
\begin{figure}[p]
    % chelsea- make full size for eprint
\begin{pchstack}[boxed,center,space=1em]
\begin{pcvstack}
\procedure[]{$\setup(1^\usecp)$} {
	(\GG, q, g) \leftarrow \GroupGen(\usecp) \\
	\pp \gets ((\GG, q, g), \HashOne, \HashThree, \HashTwo) \\
	\pcreturn \pp \\
  \text{\algcom $\pp$ is given implicitly} \\
  \text{\algcom  to all other algorithms}
}
\pcvspace
\procedure[]{$\keygen(n,t,\minsigners) $ } {
  \text{\algcom Performed by a trusted} \\
  \text{\algcom dealer, or DKG} \\
  \pcif \minsigners < 2t-1 \text{ or } \minsigners > n\\
  \pcind \pcreturn \bot \\
  \text{\algcom{Require $\geq 2t-1$ signers}} \\
  \skgroup \rgets \Zq;~\pkgroup \gets g^{\skgroup} \\
  \{(i, \sksign_i) \}_{i =1}^n \rgets \shamir.\IssueShares(\skgroup, n, t) \\
  (\skrand_i)_{i =1}^n \rgets \DVRFOnem[\HashOne].\keygen(n, t, \minsigners) \\
  \pcfor i \in \{ 1, \ldots, n \} \pcdo \\
  \pcind \pksign_i \gets g^{\sksign_i}\\
  %\mathsf{pubkeys} =  (\pksign, \{ \pksign_i \}_{i \in \set{n}})   \\
  \pcind \skparty_i \gets (\skrand_i, \sksign_i, \pkgroup) \\
  \pcind \pkparty_i \gets (\pksign_i) \\
  \pcreturn  (\pkgroup, \{ (\pkparty_i, \sk_i) \}_{i \in \set{n}}  )
}
\pcvspace
  \procedure[]{$\sign_1(\partyid, \skparty_\partyid, \msg)$} {
    (\skrand_\partyid, \sksign_\partyid, \pkgroup) \gets \skparty_\partyid \\
    y_\partyid \gets \HashThree(\pkgroup, \msg) \\
    (r_\partyid, R_\partyid) \gets \DVRFOnem[\HashOne].\generate(\partyid, \skrand_\partyid,  y_\partyid) \\
    \pcreturn (y_\partyid, R_\partyid)
}
\end{pcvstack}
  \pchspace
\begin{pcvstack}
\procedure[]{$\sign_2(\partyid, \skparty_\partyid, \msg, \coalition, \{(y_j, R_{j} )\}_{j \in \coalition})$} {
  \pcreturn \bot\ \pcif \lvert \coalition \rvert < 2t-1 \\
  (\skrand_\partyid, \sksign_\partyid, \pkgroup) \gets \skparty_\partyid \\
  y' \gets \HashThree(\pk, \msg) \\
  \pcfor j \in \coalition \pcdo \\
  \pcind \pcif y_j \neq y' \\
  \pcind \pcind \pcreturn \bot \\
  (r_\partyid, R_\partyid') \\
  \qquad \gets \DVRFOnem[\HashOne].\generate(\partyid, \skrand_\partyid, y') \\
  \text{\algcom Re-derive state from $\sign_1$} \\
  %\chelsea{$R^*$ is used as the adversarial output in the proof; please do not change this notation}
  \pcif \partyid \not\in \coalition \text{ or } R_k' \neq  R_{j}  \\
  \pcind \pcreturn \bot \\
  \mathsf{input} \gets (t, \minsigners, \coalition, \{ R_j\}_{j \in \coalition}) \\
  \pcif \DVRFOnem[\HashOne].\verify(\mathsf{input}) \neq 1 \\
  \pcind \pcreturn \bot \\
  R \gets \DVRFOnem[\HashOne].\combinepub(\mathsf{input}) \\
  c \gets \HashTwo(\pkgroup, \msg, R) \\
  z_\partyid \gets r_\partyid + c \cdot \sksign_\partyid \\
  \pcreturn z_\partyid
}
\pcvspace
  \procedure[]{$\combine(\pkgroup, \msg, \coalition, \{(y_j, R_{j}), z_j \}_{j \in \coalition}))$} {
  \mathsf{input} \gets (t, \minsigners, \coalition, \{ R_j\}_{j \in \coalition}) \\
  R \gets \DVRFOnem[\HashOne].\combinepub(\mathsf{input}) \\
  c \gets \HashTwo(\pkgroup, \msg, R) \\
  \text{\algcom $\lambda_j$ are Lagrange} \\
  \text{\algcom coefficients for $\coalition$} \\
  z \gets \sum_{j \in \coalition} z_{j} \lambda_j\\
  \pcif g^z \neq R \cdot \pk^c \\
  \pcind \pcreturn \bot \\
  \pcreturn \sigma=(R, z)
}
\end{pcvstack}
\end{pchstack}
    \caption{
      \glitter, a deterministic threshold Schnorr signature scheme.
    }
  \label{fg:construction}
\end{figure}

\section{\glitter, A Deterministic and Stateless Two-Round Threshold Schnorr Signature Scheme}

We now introduce \glitter,
an efficient,  two-round, deterministic threshold Schnorr signature scheme for moderately sized groups of participants that does not require participants to keep state between rounds of the signing protocol.
As a building block,
\glitter uses \DVRFOne to generate nonces deterministically,
and to verify that all other participants followed the protocol honestly.
\glitter is secure assuming fewer than $t$ participants are corrupted,
and at least $\minsigners$ participated in the signing protocol,
where $\minsigners \leq n$ but  $\minsigners \geq 2t-1$.

We further require that messages exchanged between participants are sent over an authenticated channel.
\glitter builds upon the verifiable pseudorandom secret sharing scheme \DVRFOne defined in Section~\ref{dvrf-concrete},
as well as Shamir's secret sharing.
Verification of signatures is identical to the single-party Schnorr verification algorithm.

\begin{remark}[Distributed Key Generation]
The \glitter construction given in Figure~\ref{fg:construction} assumes a centralized key generation procedure;
however,
using a distributed key generation (DKG) scheme is equally possible.
We discuss one possible DKG in Section~\ref{extensions}.
\end{remark}

\begin{remark}[Requirement of Authenticated Channels]
  \label{remark:auth-channel}
  We require that the messages exchanged between participants in the \glitter construction shown in Figure~\ref{fg:construction} be sent over authenticated channels;
  i.e., messages must be authenticated and verifiable as having come from their purported senders.
  Otherwise,
  an adversary can simply pick contributions that are consistent with a single honest party's $R_i$,
  which would result in a valid input for $\DVRFOnem.\verify$.
  Note that we do \emph{not} assume the authenticated channel maintains any session identifiers or state about messages;
  \glitter remains secure even if the adversary were to replay old authenticated messages.
%
  However, if a participant receives an unauthenticated message,
  we require that the participant aborts.
\end{remark}

\subsection{The Construction}

We now give more detail for each stage in \glitter;
see Figure~\ref{fg:construction} for a high-level overview.

\paragraph{\bf Key Generation.}
All participants with identifiers $i \in \set{n}$ begin by receiving a secret signing share $\sksign_i$ and a public signing share $\pksign_i = g^{\sksign_i}$.
In Figure~\ref{fg:construction},
we show key generation as a centralized procedure,
but a DKG can likewise be used.
Each $\sksign_i$ is a $t$-of-$n$ Shamir secret sharing of the group's joint secret key $\sk$;
participants use these keys for signing messages.
Participants also receive the public signing keys $\{ \pksign_i \}_{i \in \set{n}}$ for all other participants.

Each participant receives a secret \DVRFOne key $\skrand_i$ generated by performing $\DVRFOnem.\keygen$.
Participants use these keys to generate nonces and commitments for each signing session.
%
Each participant's public key share $\pkparty_i$ is just one public key,
where $\pkparty_i = \pksign_i$.
Each participant's secret key share is the tuple $\skparty_i = (\skrand_i, \sksign_i, \pkgroup)$.

\paragraph{Coordinator Role.}
Arctic assumes an external mechanism to choose the set
$\coalition \subseteq \{1, \ldots, n \}$ of signers,
such that $\minsigners \leq \abs{\coalition} \leq n$.
The coordinator may perform denial of service attacks,
but otherwise cannot impact the security of Arctic.
The coordinator can be a stand-alone entity,
or may also be a signer as well.

\paragraph{\bf Signing.}
In the first round of signing,
each participant $\partyid$ receive as input a message $\msg$.
First,
each party derives $y_\partyid \gets \HashThree(\pk, \msg)$.
To generate their nonce $r_\partyid$ and commitment $R_\partyid$,
each participant performs
$(r_\partyid, R_\partyid) \rgets \DVRFOnem.\generate(\partyid, \skrand_\partyid, y_\partyid)$.
Each participant then outputs $(y_\partyid, R_\partyid)$;
they do not need to keep any state.

In the second round of signing,
all participants again receive as input a message $\msg$,
as well as a set $\coalition$ representing the indices of at least $\minsigners$ signers,
where $\coalition \subseteq \set{n}, \lvert \coalition \rvert \geq \minsigners$.
Additionally, participants receive the list of tuples $ \{(y_j, R_j) \}_{j \in \coalition}.$
First,
each party re-derives $y' \gets \HashThree(\pkgroup, \msg)$.
Then,
each party checks the consistency of all other parties' views of $\msg$
by checking that for each $i \in \coalition$,
$y_i = y'$.
Because we require that each protocol message is authenticated by its respective party,
then this check guarantees that an adversarial player cannot split the view of honest players by sending different messages or choices of coalitions.
If any check fails,
the party aborts.

Otherwise, if all consistency checks succeed,
each participant re-derives their nonce and commitment using \DVRFOne,
again performing
$(r_\partyid, R_\partyid') \rgets \DVRFOnem.\generate(\partyid, \allowbreak \skrand_\partyid, \allowbreak y')$.
Then, the participant checks its identifier is in the coalition $C$ and that $R_\partyid' = R_j$,
aborting if either check does not hold.
Then,
each participant verifies that all other participants followed the protocol to derive  their  commitment,
by checking $\DVRFOnem.\verify(t, \minsigners, \coalition, (R_j)_{j \in \coalition})$.
If the check fails,
they abort the protocol.

Finally, if all of the above checks pass,
each participant $\partyid \in \coalition$ will then derive the group commitment $R \gets \DVRFOnem.\combinepub(t, \minsigners, \coalition, \{ R_j\}_{j \in \coalition})$,
and the challenge $c \gets \HashTwo(\pkgroup, \msg, R)$.
Finally,
each participant derives their signature share $z_\partyid \gets r_\partyid + c \cdot \sksign_\partyid$.
Each participant outputs $z_\partyid$ as its output for $\sign_2$.


\paragraph{\bf Combination and Verification.}
To perform the $\combine$ algorithm,
the group commitment $R$ is first derived using values output by participants from $\sign_1$.
Then, the response $z$ is derived by finding $z \gets \sum_{j \in \coalition} z_j \lambda_j$, where the $\lambda_j$ are the Lagrange coefficients for the set $\coalition$.
The output from $\combine$ is the Schnorr signature $\sigma = (R, z)$,
which can be verified using the single-party Schnorr verification algorithm given in Section~\ref{background:schnorr}.


\subsection{Possible Extensions}
\label{extensions}

We now discuss two possible extensions to \glitter.

\paragraph{\bf Robustness.}
\glitter as currently defined is not robust;
if any party submits invalid commitments,
then the output from $\DVRFOnem.\combinepub$ cannot be used.
%
However, it is possible to extend \DVRFOne to be robust,
therefore also ensuring that \glitter can likewise be extended.
To do so in a secure manner,
the robust extension would require additional players,
along with a protocol to come to consensus about which players misbehaved.

In particular,
by requiring that the minimum number of participants $\minparticipants$ be of size at least $\minparticipants \geq 3t-2$,
\DVRFOne and \glitter can be securely used in a robust manner.
The requirement that $\minparticipants \geq 3t-2$ is referred to as the \emph{honest supermajority} setting.
%
We give further details on how robustness can be achieved in Section~\ref{robust-shine}.
%
In this setting,
$\DVRFOnem.\verify$ can both detect any inconsistencies as well as identify the misbehaving players.
This property could likewise allow for extending \glitter to support robustness,
under the same assumption that $\minparticipants \geq 3t-2$.

\paragraph{\bf Distributed Key Generation.}
To define the signing keys $(\pksign, (\sksign_i)_{i=1}^{n})$,
any DKG that Shamir secret shares a discrete logarithm secret key could be employed,
such as the three-round DKG by Gennaro et al.\cite{GennaroJKR07}.

As one possibility to perform $\DVRFOnem[\HashOne].\keygen$ in a distributed manner,
each subset of $n-t+1$ parties could assign one representative to generate a secret PRF key uniformly at random for that subset.
The representative would then share this PRF key with all parties in their subset.
To ensure security against fully malicious adversaries,
parties within each subset could perform a subsequent broadcast round to compare their received inputs,
to ensure consistency.

\subsection{Security}
\label{arctic-security}

\paragraph{\bf Correctness.}
In the first round of signing,
participants will output nonces and commitments
$(r_i, R_i) \gets \DVRFOnem.\generate(i, \skrand_i, y_i)$,
for $i \in \coalition$,
where $y_i \gets \HashThree(\pkgroup, \msg)$,
and $R_i = g^{r_i}$.

In the second round of signing,
each participant receives the set of tuples $\{(y_j, R_j) \}_{j \in \coalition}$.
After deriving $y' \gets \HashThree(\pkgroup, \msg)$,
then the check that $y_j = y'$ for each $j \in \coalition$ will succeed.

Because \DVRFOne is correct,
then
$\DVRFOnem.\verify(t, \minsigners, \coalition, \{R_j\}_{j \in \coalition})$ will output $1$
and the same group commitment $R = \prod_{j \in \coalition} R_j^{\lambda_j}$ will be computed regardless of the choice of $\coalition$.
Because $\DVRFOnem.\generate$ is deterministic,
then after deriving $y \gets \HashThree(\pkgroup, \msg)$,
$\DVRFOnem.\generate(\partyid, \allowbreak \skrand_\partyid, y)$
will output the same $(r_\partyid, R_\partyid)$ as derived in the first round of signing.


After deriving $(r_\partyid, R_\partyid)$,
all signers in a coalition $\coalition$ output valid signature shares $z_\partyid$
with respect to $R$
and challenge $c = \HashTwo(\pkgroup, m, R)$,
where $z_\partyid =  r_\partyid + c \cdot \sksign_\partyid$.
The aggregated signature is then $\sigma = (R, z)$,
where $z = \sum_{j \in \coalition} z_j \cdot \lambda_j$.
Because $\skgroup = \sum_{j \in \coalition} \sksign_j \lambda_j$,
$\pkgroup = g^\skgroup = g^{\sum_{j \in \coalition} \sksign_j \lambda_j}$,
and $R = \prod_{j \in \coalition} R_j^{\lambda_j} = g^{\sum_{j \in \coalition} r_j\cdot \lambda_j}$,
we have that $g^z = R \cdot \pkgroup^c$, as required.


